\section{Introduction}
\label{sec:tpc33-introduction}

The purpose of this paper is to turn one exceptional Fourier coefficient
into a concrete arithmetic average.  TPC-32 proves determinant-frequency
dispersion for a matched prime--M\"obius cutoff shell, but the physical
problem selects the coefficient at frequency zero
\citep{WangTPC32}.  Ordinary frequencies are small outside a power-small
exceptional set; membership of the distinguished coefficient in that set
is not decided by Parseval.  The same paper shows that support size,
\(\ell^1\)-mass, \(\ell^2\)-energy, and even perfect cancellation at every
nonzero frequency do not force cancellation at zero.  Additional
arithmetic structure is required.

There are two logically separate questions.  First, what is the exact
arithmetic sequence inside a zero-frequency column after the actual
row-pair mask is restored?  Second, can averaging over the physical family
of those sequences replace a pointwise two-point M\"obius theorem?  The
present paper answers the first question and proves a finite-field model
for the second.  It does not complete the final transfer to the integer
physical packet.

\subsection{From a bounded mask to the actual column}

The abstract interface of TPC-25--TPC-32 deliberately retains the row-pair
mask as an arbitrary bounded joint multiplier
\citep{WangTPC25,WangTPC32}.  That is the correct interface for the earlier
norm estimates, but it is too broad for an affine-column theorem: an
arbitrary entrywise bounded matrix can align with any prescribed kernel.
We therefore return to the actual mask inherited from TPC-18 and recorded
in TPC-21 \citep{WangTPC21},
\begin{equation}
 \mathfrak m((\ell,d),(\ell',e))
 =\one_{\ell\ne\ell'}
  \one_{|\ell d-\ell'e|>QX^{-\kappa_0}}
  \one_{(d,e)\le X^{\kappa_0}}.
 \label{eq:tpc33-intro-actual-mask}
\end{equation}
Here \(d,e\) are squarefree and lie on one opened dyadic scale.  For a
fixed opposite row \((\ell',e)\), the last indicator has a divisor
expansion of total coefficient mass \(3^{\omega(e)}=X^{o(1)}\), while the
middle indicator merely removes one interval from the \(d\)-line.  After
one ultra divisor is reflected, the remaining inner variable lies on
\begin{equation}
 \widetilde d_r=\widetilde d_0+sr,
 \qquad
 \widetilde u_r=\widetilde u_0+\ell jv r,
 \qquad
 s\widetilde u_0-\ell jv\widetilde d_0=h.
 \label{eq:tpc33-intro-affine-column}
\end{equation}
Thus the actual mask has small \emph{columnwise} expansion complexity.
This does not say that the entire mask matrix has small projective or Schur
norm.

\subsection{Exact arithmetic of a fixed-determinant line}

The equation
\begin{equation}
 sU-aD=h
 \label{eq:tpc33-intro-determinant-line}
\end{equation}
is not merely a nonproportional pair of affine forms.  Under the physical
primitivity conditions \((a,s)=1\) and \((as,h)=1\), it belongs to one
unimodular orbit.  Choosing \(sx-ay=1\) gives
\[
 \binom{U_t}{D_t}
 =\begin{pmatrix}x&a\\y&s\end{pmatrix}
  \binom{h}{t}.
\]
The determinant-one matrix transfers all common divisors back to the two
coordinates \((h,t)\).  In particular, the primitive part of the two-sign
sequence is one M\"obius value on the reducible quadratic
\begin{equation}
 D_tU_t
 =as t^2+h(sx+ay)t+h^2xy,
 \qquad
 \operatorname{disc}(D_tU_t)=h^2.
 \label{eq:tpc33-intro-quadratic}
\end{equation}
This reformulation does not prove signed cancellation.  We do prove a
uniform squarefree-support theorem for a balanced B\'ezout origin.  Its
main term has a positive Euler product uniformly in the growing slopes
when \(\max(a,s)=o(T)\) by a fixed power.

\subsection{Projective averaging at fixed determinant}

The central unconditional average theorem is finite and exact.  Over a
prime field, the change of variables
\begin{equation}
 (a,s)\longmapsto(b,c)=(a/s,h/s)
 \label{eq:tpc33-intro-projective-change}
\end{equation}
is a bijection of \((\F_q^\times)^2\).  Translation Parseval in \(c\),
followed by summation over the slope \(b\), gives an exact second moment for
every fixed nonzero determinant \(h\).  Outside \(O(q^{2-2\eta})\)
coefficient pairs, it yields an \(O(q^{1/2+\eta})\) centered fluctuation and
hence \(o(q)\) total correlation in the finite M\"obius model.
Moreover, on a complete residue orbit the restriction \(a=\ell j\) loses
no energy: multiplication by every fixed nonzero \(\ell\), prime or not,
permutes \(j\in\F_q^\times\).

This theorem rotates the earlier ``almost every determinant frequency''
statement into ``fixed determinant, almost every nonzero coefficient pair.''
The average runs over both \(b=a/s\) and the affine offset \(c=h/s\); by
itself it is not an almost-every-\(b\) statement.  A
direct physical consequence still fails for a precise reason.  Integer
no-wrap needs a modulus on the ambient determinant scale, approximately
\(X\), whereas complete coverage of the orbit parameter is available only
for moduli at most \(J=X^{133/400+o(1)}\).  A single modulus cannot satisfy
both requirements.

\subsection{The exact conditional gate}

Let \(C_m\) denote the ultra increment and \(\mathsf H_m\) the half-prefix
from the exact polarization in TPC-32.  The full matched shell is
\begin{equation}
 \cK^{\rm sh}_{\alpha,\gamma}(j)
 =C_{m_\alpha}(j)\mathsf H_{m_\gamma}(j)
  +\mathsf H_{m_\alpha}(j)C_{m_\gamma}(j).
 \label{eq:tpc33-intro-polarization}
\end{equation}
Keeping the opposite row outside each column defines two nonnegative
column energies \(E_L,E_R\).  The inherited one-row estimate
\(\|\gamma^{(i)}\|_2^2\ll Q\log(2L)\) gives the rigorous implication
\begin{equation}
 E_L+E_R\ll X^{o(1)}Q^3
 \quad\Longrightarrow\quad
 |\cS_{\rm sh}^{\rm all}|\ll X^{o(1)}Q^2.
 \label{eq:tpc33-intro-energy-transfer}
\end{equation}
The left side is a new analytic input, not a consequence of the row-energy
bound.  Separately estimating the two polarizations is sufficient but
possibly stronger than estimating their scalar sum.

At the source-compatible endpoint of TPC-28,
\begin{equation}
 Q=X^{267/400+o(1)},\qquad
 J=X^{133/400+o(1)},\qquad
 JQ\asymp X,
 \label{eq:tpc33-intro-high-beta}
\end{equation}
and therefore \(Q^2=(JQ^2)J^{-1}\).  The right side of
\eqref{eq:tpc33-intro-energy-transfer} is exactly the amplitude budget of
the preceding zero-frequency gate.  It is important not to restate this as
a proved flatness ratio: no matching lower bound for the determinant-cell
energy is known.

\subsection{Result ledger and scope}

The logical status of the main statements is summarized in
\cref{tab:tpc33-result-ledger}.

\begin{table}[t]
\centering
\caption{Logical status of the TPC-33 interfaces.}
\label{tab:tpc33-result-ledger}
\begin{tabular}{@{}>{\raggedright\arraybackslash}p{0.39\textwidth}
>{\raggedright\arraybackslash}p{0.16\textwidth}
>{\raggedright\arraybackslash}p{0.35\textwidth}@{}}
\toprule
Statement & Status & What it controls \\
\midrule
\(SL_2(\Z)\) affine normal form and gcd identities
& unconditional & every primitive line \(sU-aD=h\) \\
Balanced-origin squarefree energy
& unconditional & support, not signs \\
Actual-mask divisor expansion
& unconditional & one fixed physical column \\
Fixed-\(h\) projective Plancherel
& unconditional & complete two-parameter finite family \\
Finite M\"obius almost-all-coefficient-pair bound
& unconditional & modular model, not integer shell \\
\(Q^3\) column energy \(\Rightarrow Q^2\) shell
& conditional transfer & distinguished physical zero mode \\
Restricted physical \(Q^3\) energy
& open & next analytic gate \\
\bottomrule
\end{tabular}
\end{table}

The new contribution should be read as an exact synthesis tailored to the
TPC residual: the structured physical disintegration, the fixed-determinant
projective identity with its prime-product lift, and the quantified energy
splice.  Basic finite Fourier orthogonality and squarefree sieving have a
large classical literature, so no priority claim is made for those
ingredients in isolation.  Nor do the results imply a Hardy--Littlewood
main term, positivity of a prime-pair correlation, infinitely many twin
primes, or a general resolution of the sieve parity problem.

The paper is organized as follows.  \Cref{sec:tpc33-physical-interface}
fixes the matched shell and the actual mask.  \Cref{sec:tpc33-affine-normal}
proves the integral normal form, and \cref{sec:tpc33-squarefree-energy}
proves the support theorem.  \Cref{sec:tpc33-structured-mask} performs the
physical column disintegration.  \Cref{sec:tpc33-projective-plancherel}
proves the finite-field identity, and
\cref{sec:tpc33-mobius-corollaries} derives its M\"obius and product-slope
consequences.  \Cref{sec:tpc33-column-transfer} gives the \(Q^3\) transfer.
The final sections record sharp obstructions, the literature boundary, the
high-\(\beta\) ledger, and the exact computational certificate.
